\documentclass{article}

% NeurIPS 2026 style. "preprint" = non-anonymous, no line numbers ("final" is reserved for accepted NeurIPS papers).
\usepackage[preprint]{neurips_2026}

\usepackage[utf8]{inputenc}
\usepackage[T1]{fontenc}
\renewcommand{\ttdefault}{lmtt}
\usepackage{hyperref}
\usepackage{url}
\usepackage{booktabs}
\usepackage{amsfonts,amsmath}
\usepackage{microtype}
\usepackage{xcolor}
\usepackage{graphicx}
\usepackage{float}

\newcommand{\sone}{\textsc{S1}}
\newcommand{\stwo}{\textsc{S2}}
\newcommand{\todo}[1]{\textcolor{red}{[TODO: #1]}}

\title{Reading Robustness from One Truncated Trace:\\Per-Model White-Box Recipes for Spurious Math Reasoning}

\author{%
  \todo{Author 1} \\ \todo{Affiliation} \\ \texttt{abhichoudhury19@gmail.com} \And
  \todo{Author 2} \\ \todo{Affiliation} \\ \todo{email}
}

\begin{document}
\maketitle

\begin{abstract}
The AIMO Interpretability Challenge asks whether a reasoning model's success on a competition math problem survives faithful rewordings of the problem.
We predict this from a single greedy reasoning trace truncated at 6144 tokens, with white-box access to the model.
Our method (i) calls every problem whose reasoning finishes within the budget robust, and (ii) scores unfinished traces with per-model, label-free-calibrated readouts: a rank ensemble of self-verification share, answer survival when the reasoning is cut, hidden-state sensitivity to automatic rewordings, attention look-back and forced-answer confidence; for Qwen3.5-4B, the entropy of the forced answer after 4096 tokens; for Olmo-3-7B-Think, a Min-K\% familiarity router and a second model acting as a judge.
Our two submissions differ only in Olmo's unfinished calls: the second replaces them with an early-trace hidden-state direction.
They score 1.00 and 0.93 on the public leaderboard (all-robust: 0.50) and make 33 and 36 correct calls on the official training labels for our four models (43 rows; all-robust: 31). Removing the hidden-state and attention signals lowers leaderboard accuracy to 0.86, and removing the output-distribution signals instead lowers it to 0.71--0.79.
We also report a negative result: labels we generated with our own perturbations do not reproduce the official labels, which shaped how we validated everything.
Code: \todo{GitHub URL}.
\end{abstract}

\section{Introduction}
\label{sec:intro}
Both nominated submissions share one test-time pipeline (Section~\ref{sec:method}) and differ only in how unfinished Olmo problems are called, so this single report covers both, as the rules require.
\sone{} (\texttt{permodel-v8-ent4096}) is the base method. \stwo{} (\texttt{permodel-v8-ent4096-olmoseg25}) is \sone{} with one single-feature \emph{override} for Olmo (Section~\ref{sec:overrides}).
One finding shaped the design: robustness labels that we generated ourselves did not reproduce the official ones (Section~\ref{sec:labels}).

\textbf{Contributions.} (1) A test-time pipeline that needs one greedy trace plus a few teacher-forced passes per problem. (2) Per-model feature recipes with label-free thresholds. (3) A component ablation on the official labels.

\section{Data and evaluation}
\label{sec:data}
\textbf{Models.} Qwen3.5-4B \citep{qwen35}, Skywork-OR1-Math-7B \citep{he2025skywork}, Olmo-3-7B-Think \citep{olmo3} and DeepSeek-R1-0528-Qwen3-8B \citep{deepseek2025r1,yang2025qwen3}; perturbation robustness in the sense of \citet{mirzadeh2025gsmsymbolic}, labelled by the challenge \citep{stefanik2026aimointerp}.

\textbf{Official labels.} \textsc{Main}: the organisers' training set \texttt{train-main-v2}, restricted to our four models: 43 rows (Qwen 12, Skywork 12, Olmo 9, DeepSeek 10), 31 robust. This is the only labelled data we evaluate on; we report correct calls out of 43 with Wilson 95\% intervals \citep{wilson1927}.
\textbf{Public leaderboard.} We report leaderboard accuracy on its 14 problems, and also use it as a held-out validation set. However, with only 3--5 problems per model, estimates from this set are very noisy. The problems and the evaluation did not change during the competition, so builds are scored from their logged calls where they were not submitted separately.

\textbf{Our own data, used for calibration only.} (a) Unlabelled greedy traces for $\sim$100 post-cutoff competition problems per model (\textsc{Pool}). They supply the reference distributions and median thresholds below, which need no labels. (b) Problems we labelled ourselves. No official threshold for the accuracy drop is published; we inferred the rule from the data. In \textsc{Main}, every robust row has a worst-case drop $\le 0.10$ and every non-robust row $\ge 0.25$, with nothing in between, and the same rule, counting never-solved problems as robust, reproduces the labels of the public validation sample exactly. We used 8 temperature-1 samples per version and five perturbation types written by an LLM. Every generation of our own was capped at 6144 new tokens, matching the test-time trace, with answers forced at the cap when the model had not finished (the relabel in Section~\ref{sec:labels} used 32k); the official labels come from a much larger generation budget. These labels set feature \emph{signs} and, for \stwo{}, fit one direction.
Section~\ref{sec:labels} shows that these labels are unreliable, so no result below is measured on them.

\textbf{Regime split.} At a 6144-token cap most \textsc{Main} traces are still thinking (32/43).
All 11 \textsc{Main} rows that finish are robust. Every non-robust label therefore falls on an \emph{unfinished} trace, which is where every recipe acts.

\section{Method}
\label{sec:method}
\textbf{Trace.} We build the challenge prompt (``Solve the following mathematical problem step by step. Put your final answer within \texttt{\textbackslash boxed\{\}}'') with each model's chat template. We decode greedily (bf16, HF Transformers \citep{wolf2020transformers}) up to 6144 new tokens, stopping early on end-of-thinking, end-of-sequence or a 3$\times$32-token loop.
\textbf{Finished} (the end-of-thinking token appears) $\Rightarrow$ robust. \textbf{Unfinished} $\Rightarrow$ score with the model's recipe.

All features are computed on that one trace (definitions in Appendix~\ref{app:features}). We group them by what they read; Section~\ref{sec:wb} ablates the groups.

\subsection{Text features (black-box)}
\textbf{Finished rule}: whether the end-of-thinking token appears within the cap. \texttt{chk}: the share of $\ge$4-word sentences of the reasoning that check, verify or confirm something.

\subsection{Output-distribution features (grey-box)}
These read the model's next-token distribution under teacher forcing.
\texttt{fa}: the top-1 probability of the first answer token after appending ``\texttt{</think>} The final answer is \texttt{\$\textbackslash boxed\{}'' at 6144 tokens.
$H_{b}$: the entropy of that distribution when the trace is cut at $b$ tokens.
\emph{Answer survival} \citep[cf.\ early answering,][]{lanham2023faithfulness}: \texttt{nrk} and \texttt{f1k}, the probability of the clean forced answer when the reasoning is removed entirely or cut to its first 1024 tokens, relative to the clean probability.
\texttt{cfl}: the change in the trace's log-likelihood when it is teacher-forced after three automatic rewordings of the problem (10\% typos, a neutral distractor sentence, renamed variables).
Min-K\% \citep{shi2024mink} ($k=20$): the mean of the lowest 20\% token log-probabilities of the problem text, a familiarity score.

\subsection{Internal features (white-box)}
These read activations directly.
\emph{Hidden states}: \texttt{cfc}, the cosine between late-trace hidden states of the original pass and of the three reworded passes; the Olmo direction of \stwo{} (Section~\ref{sec:overrides}).
\emph{Attention}: \texttt{lble} and \texttt{pmle}, the Lookback-Lens ratio \citep{chuang2024lookback} and the attention mass on the prompt, each as last quarter minus first quarter of the trace; \texttt{jmp}, the attention-weighted mean backward distance between 64-token blocks (robust traces keep re-reading their own setup).
The Qwen judge used for Olmo (below) is also internal, since Qwen's ensemble consists of \texttt{lble}, \texttt{pmle}, \texttt{jmp} and \texttt{cfc}.

\subsection{From features to calls}
\textbf{Rank-ensemble score.} For model $m$ with feature subset $F_m$, a value $x_f$ is mapped to its percentile $q_f$ among the reference values $R_{m,f}$ (\textsc{Pool} traces plus our labelled traces, unfinished only):
\begin{equation}
  s = \frac{1}{|F_m|}\sum_{f\in F_m}\big(\tilde q_f - \tfrac12\big),\qquad \tilde q_f = \begin{cases} q_f & \sigma_f=+1\\ 1-q_f & \sigma_f=-1\end{cases},\qquad \text{robust} \iff s>\tau_m .
\end{equation}
$\tau_m$ is the median score of model $m$'s calibration traces, which needs no labels. The signs $\sigma_f$ are the direction of each feature's length-controlled AUC.
The subsets $F_m$ were chosen per model by subset search on our labelled rows:
Qwen \{\texttt{lble, pmle, jmp, cfc}\} (Qwen's own unfinished calls use $H_{4096}$ instead, below; this ensemble serves as the judge for Olmo); Skywork \{\texttt{fa, chk, f1k}\}; DeepSeek \{\texttt{chk, nrk, cfl, pmle}\}.
Percentile scoring makes the score insensitive to feature scale and to shifts between our hardware and the evaluator's.

\textbf{Olmo: novelty router and judge.} Olmo's features flipped direction between problems it had seen in training and post-cutoff problems.
We route on Min-K\%: below $-5.93$ (the cut that best separated post-cutoff from older problems in \textsc{Pool}, using no robustness labels) goes to a \emph{novel} branch \{\texttt{f1k}\}, otherwise to a \emph{familiar} branch \{\texttt{chk, lble, pmle}\}, each with its own median threshold.
A \emph{judge} then loads Qwen, generates Qwen's own trace of the same problem and scores it with Qwen's recipe. Olmo's call is robust iff the mean of two label-free percentiles exceeds $0.5$: Olmo's margin among its \textsc{Pool} margins, and Qwen's margin among Qwen's margins on the same \textsc{Pool} problems. If the judge fails or time is short, Olmo's own call stands.

\textbf{Qwen: forced-answer entropy at 4096 tokens.} For unfinished Qwen problems, \sone{} replaces the rank-ensemble call with $H_{4096}$ (applied last, after any judging): robust iff $H_{4096} < 1.621$, the median over Qwen's 100 \textsc{Pool} traces. The sign comes from our 66 labelled Qwen problems (AUC $0.77$ for lower-is-robust). Qwen's rank-ensemble margins are still computed, because the Olmo judge uses them.

\subsection{\stwo: an early-trace hidden-state direction for Olmo}
\label{sec:overrides}
\stwo{} runs the full \sone{} pipeline and then, \emph{after} the judge, replaces the call for unfinished Olmo problems. Every other model and every finished problem keeps its \sone{} call.
We take the mean hidden state over the 256 tokens around 25\% of the trace, at the 5th decoder layer, and project it onto a unit robust-minus-non-robust difference-of-means direction fitted on our 21 labelled unfinished Olmo problems (shipped with the submission). Robust iff the projection is above $0.234$, the median over Olmo's 50 \textsc{Pool} traces.
For both single-feature readouts the threshold is label-free and the sign comes from our data. Choosing \emph{which} feature to use for these two models, however, used \textsc{Main} results and leaderboard scores of earlier builds (see Limitations).

\textbf{Cost.} Per unfinished problem: one generation of up to 6144 tokens, $\sim$10 teacher-forced passes (forced answers at several cuts, three rewordings, attention hooks) and, for Olmo, one Qwen generation. Everything runs on a single GPU within the challenge time limit (one problem per call).

\section{Results}
\label{sec:results}
\begin{table}[t]
\caption{Correct calls on \textsc{Main} (43 rows) with Wilson 95\% intervals, and public leaderboard accuracy (14 problems). On \textsc{Main}, Olmo's judge cannot be replayed (it needs a Qwen trace per Olmo problem), so Olmo uses its own branch there. \sone{} is best on the leaderboard, \stwo{} on \textsc{Main}; on \textsc{Main} both gains over all-robust are inside the intervals.}
\label{tab:main}
\centering\small
\begin{tabular}{lcc}
\toprule
Method & \textsc{Main} correct (43) & Leaderboard score \\
\midrule
All robust (baseline) & 31 \; [0.57, 0.83] & 0.50 \\
Finished $\Rightarrow$ R, unfinished $\Rightarrow$ NR & 23 \; [0.39, 0.67] & 0.57 \\
\textbf{\sone{}} (\texttt{v8-ent4096}) & 33 \; [0.62, 0.87] & \textbf{1.00} \\
\textbf{\stwo{}} (\sone{} + Olmo direction) & \textbf{36} \; [0.70, 0.92] & 0.93 \\
\bottomrule
\end{tabular}
\end{table}

Table~\ref{tab:main} gives the headline numbers. On the leaderboard \sone{} scores 1.00 and \stwo{} 0.93, against 0.50 for all-robust. \textsc{Main} is 72\% robust, so calling everything robust is a strong baseline there; \stwo{} beats it by 5 calls and \sone{} by 2, within the intervals. The two submissions trade places across the two label sets, which is why we nominate both. Among all builds we scored on the leaderboard (Appendix~\ref{app:ranking}), the per-model builds scored 0.64--1.00 and single-family builds 0.28--0.71.

\section{Ablations}
\label{sec:ablations}
\begin{table}[t]
\caption{Component ablation: correct calls per model on \textsc{Main} (Qwen 12, Skywork 12, Olmo 9, DeepSeek 10 rows; of them 11/7/4/10 unfinished) and leaderboard accuracy (LB). The per-model recipes help DeepSeek and Skywork most; the Qwen entropy helps on both sets; the Olmo direction helps on \textsc{Main} but costs one leaderboard problem.}
\label{tab:ablation}
\centering\small
\begin{tabular}{lcccccc}
\toprule
Variant & Qwen & Skywork & Olmo & DeepSeek & Total & LB \\
\midrule
All robust                         & 8 & 9  & 7 & 7 & 31 & 0.50 \\
Finished only (unfinished NR)      & 5 & 8  & 7 & 3 & 23 & 0.57 \\
v8 (rank ensembles, no overrides)  & 8 & 10 & 5 & 9 & 32 & 0.93 \\
v8 without Olmo router/judge       & 8 & 10 & 6 & 9 & 33 & -- \\
\sone{} = v8 + Qwen $H_{4096}$     & 9 & 10 & 5 & 9 & 33 & 1.00 \\
v8 + Olmo direction only           & 8 & 10 & 8 & 9 & 35 & 0.86 \\
\stwo{} = \sone{} + Olmo direction & 9 & 10 & 8 & 9 & 36 & 0.93 \\
\bottomrule
\end{tabular}
\end{table}

\textbf{Components (Table~\ref{tab:ablation}).}
(i) The regime split alone is not a predictor (``finished only'' is below all-robust), but it settles 11/43 \textsc{Main} rows at zero error.
(ii) The per-model rank ensembles (v8) help most for DeepSeek (+2 over all-robust) and Skywork (+1); for Qwen they tie all-robust.
(iii) Replacing Qwen's ensemble with $H_{4096}$ (\sone{}) adds 1 \textsc{Main} call and one leaderboard problem.
(iv) The Olmo direction (\stwo{}) adds 3 \textsc{Main} calls but costs one leaderboard problem. On \textsc{Main} the Olmo router alone is $-1$ against the plain Olmo recipe; the judge, which cannot be replayed on \textsc{Main}, is not measured there. \sone{} keeps the router and judge for Olmo and is the hedge against the Olmo direction, which rests on only 4 unfinished \textsc{Main} rows.

\subsection{Does the white-box signal matter?}
\label{sec:wb}
\begin{table}[t]
\caption{White-box ablation by feature group. \emph{No internals}: every hidden-state and attention feature, the Olmo direction and the Qwen judge removed; the two submissions then coincide, since they differ only in an internal readout. \emph{No output}: \texttt{fa}, \texttt{nrk}, \texttt{f1k}, \texttt{cfl}, the Min-K\% router and $H_{4096}$ removed (Qwen falls back to its internal ensemble). \emph{Text only}: finished rule + \texttt{chk} for every model. A model left with no feature falls back to \texttt{chk}. Modified recipes use threshold 0, the label-free expected median of a percentile score (\sone{} gives 32 / 1.00 at threshold 0). Neither group alone reaches the full submissions on the leaderboard; together they lift it from 0.86 (output) or 0.71--0.79 (internal) to 1.00 and 0.93.}
\label{tab:wb}
\centering\small
\begin{tabular}{lcccccc}
\toprule
Variant & Qwen & Skywork & Olmo & DeepSeek & \textsc{Main} & LB \\
\midrule
All robust & 8 & 9 & 7 & 7 & 31 & 0.50 \\
Text only & 10 & 9 & 8 & 7 & 34 & 0.57 \\
Text + output (no internals; \sone{} = \stwo{}) & 9 & 10 & 8 & 8 & 35 & 0.86 \\
\sone{}: text + internal (no output) & 8 & 9 & 6 & 7 & 30 & 0.71 \\
\stwo{}: text + internal (no output) & 8 & 9 & 8 & 7 & 32 & 0.79 \\
\sone{} (all three) & 9 & 10 & 5 & 9 & 33 & 1.00 \\
\stwo{} (all three) & 9 & 10 & 8 & 9 & 36 & 0.93 \\
\bottomrule
\end{tabular}
\end{table}

Table~\ref{tab:wb} removes each feature group in turn. In \sone{} the internal signals are Olmo's attention-based familiar branch and Qwen judge and DeepSeek's prompt-attention feature; \stwo{} adds the Olmo hidden-state direction. Qwen's own calls come from the output distribution ($H_{4096}$) in both.
\textbf{Removing internals} lowers leaderboard accuracy from 1.00 to 0.86 for \sone{} and from 0.93 to 0.86 for \stwo{}, through the Olmo problems, where a check-share call replaces the internal readouts. On \textsc{Main} the variant without internals scores 35/43, between \sone{} (33) and \stwo{} (36).
\textbf{Removing output features} hurts more: 0.71 and 0.79 on the leaderboard and 30--32/43 on \textsc{Main}, at or below all-robust there. The internal readouts are not a substitute for the forced-answer features; they add to them.
\textbf{Removing both} (text only) leaves 0.57 and 34/43.
Over both sets together, \stwo{} with all three groups is best (49/57, against 47 without internals, 43 without output and 42 text only); \sone{} ties its no-internals variant (47/57) because it gains on the leaderboard what it loses on \textsc{Main}. \textsc{Main} is 72\% robust and was partly used for feature signs, so it favours recipes that call more problems robust; the leaderboard is balanced.

\begin{table}[t]
\caption{Where in the trace to read the signal: AUC on unfinished traces (robust = positive class, oriented by the shipped sign) on our labels / \textsc{Main}. Left: Qwen forced-answer entropy at truncation $b$ (\textsc{Main}: 7 R / 4 NR). Right: Olmo direction projections by trace position (\textsc{Main}: 2 R / 2 NR). Olmo's ``ours'' AUCs are in-sample, since the direction was fitted on those labels. On our labels the entropy signal grows with $b$; on \textsc{Main} every budget is similar, so the exact budget matters little.}
\label{tab:position}
\centering\small
\begin{tabular}{lcc}
\toprule
Qwen $H_b$, $b=$ & ours & \textsc{Main} \\
\midrule
512  & 0.44 & 0.75 \\
1024 & 0.49 & 0.64 \\
2048 & 0.61 & 0.71 \\
3072 & 0.69 & 0.75 \\
\textbf{4096} & 0.77 & 0.71 \\
6144 & 0.79 & -- \\
\bottomrule
\end{tabular}\hfill
\begin{tabular}{lcc}
\toprule
Olmo direction at & ours & \textsc{Main} \\
\midrule
\textbf{25\% of trace} & 0.99 & 1.00 \\
50\% of trace & 0.97 & 1.00 \\
75\% of trace & 0.93 & 1.00 \\
\bottomrule
\end{tabular}
\end{table}

\textbf{Truncation point and position (Table~\ref{tab:position}).} On our labels the forced-answer entropy is uninformative before about 2048 tokens and rises to 0.77--0.79 at 4096--6144 tokens: answer confidence becomes informative only after a few thousand tokens of reasoning. On \textsc{Main}, with 11 unfinished Qwen rows, all budgets give 0.64--0.75. We ship 4096, chosen with \textsc{Main} results and the leaderboard scores of earlier builds in view (see Limitations).
For Olmo, the direction separates \textsc{Main}'s four unfinished rows at every position; with so few rows the choice of 25\% is not supported over 50\% or 75\%. This is the weakest-supported component, and the reason \sone{} is nominated alongside \stwo{}.

\textbf{Alternatives tried (Appendix~\ref{app:ranking}).} We also tested whole-trace hidden-state geometry (K-means walk graphs, $\delta$-hyperbolicity and Poincar\'e embeddings \citep{raj2026hyperbolic}, CKA \citep{kornblith2019cka}), attention-map spectra \citep{binkowski2025lapeig}, Min-K\%++ \citep{zhang2025minkpp}, kNN and TDA trajectory features, PEARL-style prompt noising, and sentence-embedding reasoning graphs. Several reach high AUC on one label set but none holds up across our labels and \textsc{Main}; builds based on single families scored 0.28--0.71 on the leaderboard.

\section{What do the official labels measure?}
\label{sec:labels}
Throughout the competition we tried to enlarge the labelled data by running the official rule ourselves. Two findings made us stop trusting those labels.
\textbf{(1) Noise.} With 8 samples per version, a problem with mid-range accuracy crosses the $0.25$ drop threshold by chance alone. Our first labelled problems were harder for the models than \textsc{Main}'s: on unfinished traces the median forced-answer top-1 probability at 6144 tokens was 0.54 (ours) vs 1.00 (\textsc{Main}) for Qwen and 0.92 vs 1.00 for Skywork. We therefore labelled a new, \emph{confidence-matched} set: unfinished traces whose forced top-1 probability is at least 0.7, as for most unfinished \textsc{Main} traces. On these 109 problems the number of non-robust labels we obtained (10/10/8 for Qwen/Olmo/Skywork) matched the number expected under a no-effect null (9.3/10.3/7.1).
\textbf{(2) Non-reproduction.} \textsc{Main} problem \texttt{7ee9d443} is officially non-robust for all four models (drops 0.40--1.00). Under our perturbations Qwen and Skywork still solve it 47--48/48 times at a 32k cap, and 67--68/68 times across 30 further independent rewordings, with changed numbers, and under bare and R1-style prompts. The sampling settings match the official ones, so the official non-robust labels depend on perturbation texts or grading details we do not have.
On a 32k natural-answer relabel of 62 Skywork problems, non-robust cases were rare (8) and every recipe, including our Skywork recipe (shared by \sone{} and \stwo{}), was at chance in balanced accuracy (0.52). In the same data, original-problem accuracy predicted the label with AUC 0.74: our labels track difficulty.
We therefore evaluate only on \textsc{Main} and the leaderboard, and used our own labels only for signs.

% ---------------------------------------------------------------- not counted towards the 5 pages
\section*{Limitations}
\label{sec:limits}
\textbf{Small, partly in-sample evaluation.} \textsc{Main}'s 43 rows give intervals about $\pm 0.12$ wide; \sone{} and \stwo{} differ by 3 \textsc{Main} calls and one leaderboard problem, neither of which is significant.
\textsc{Main} is not fully held out: feature signs for the shared features used \textsc{Main} (length-controlled AUC), reference sets include \textsc{Main}'s unfinished values, and both single-feature readouts, as well as the per-model subsets of later builds, were chosen using \textsc{Main} results and the leaderboard scores of earlier builds.
\textbf{Unreliable self-generated labels} (Section~\ref{sec:labels}) leave the signs and the Olmo direction resting on 21--66 noisy labels per model.
\textbf{Distribution shift.} Features calibrated on our data transferred unevenly to leaderboard scores (single-family builds 0.28--0.71), so accuracy on the private 43 problems may differ substantially from \textsc{Main}.
\textbf{Replay gap.} Olmo's judge was not replayed on \textsc{Main}.
\textbf{Greedy-only view.} We observe one greedy trace capped at 6144 tokens, while the labels come from temperature-1 samples with a much larger budget.

\section*{AI usage statement}
\todo{Check and edit.} We used AI assistants extensively.
\textbf{Claude Code, mainly with Claude Opus 5.5 (Anthropic),} was used for much of the code and experimentation: it wrote most of the analysis, feature-extraction, labelling and submission code under our direction; ran and monitored the GPU experiments; produced the ablation tables; and drafted this report, which we then edited.
\textbf{Claude Sonnet} (through claude.ai) generated the perturbation texts for our own labelled data, following a written specification (Section~\ref{sec:data}).
\todo{List any other model or tool used, e.g.\ for literature search.}
The authors chose the research directions, decided every submission and nomination, checked the reported numbers against the experiment logs, and are responsible for all claims. No AI system is an author.

\bibliographystyle{plainnat}
\bibliography{refs}

\appendix
\section{Feature definitions}
\label{app:features}
All features are computed on the prompt $p$ and the greedy trace $t$ (at most 6144 tokens), with suffix $\sigma$ = ``\texttt{\textbackslash n</think>\textbackslash n\textbackslash nThe final answer is \$\textbackslash boxed\{}''.
\begin{itemize}
\item \texttt{chk}: share of the thinking's sentence units ($\ge$4 words) that match a check/verify/confirm lexicon; undefined below 10 units.
\item \texttt{fa}: $\max_v P(v \mid p, t_{1:6144}, \sigma)$.
\item \texttt{nrk}, \texttt{f1k}: with $a^*$ the clean forced answer token, $P(a^*\mid p,\sigma)/P(a^*\mid p,t,\sigma)$ and $P(a^*\mid p,t_{1:1024},\sigma)/P(a^*\mid p,t,\sigma)$.
\item \texttt{cfc}, \texttt{cfl}: for each automatic rewording $p'$ (typos at rate 0.1, a fixed distractor sentence after the first sentence, single-letter variables in math renamed consistently), teacher-force $t$ after $p'$. \texttt{cfc} is the mean cosine between original and reworded 16-token-window hidden states over the last quarter (averaged over 6 layers); \texttt{cfl} is the mean change in per-token log-likelihood of $t$ over its first tokens.
\item \texttt{lble}, \texttt{pmle}: per 16-token span, all layers and heads, the Lookback-Lens ratio (attention on the prompt over attention on prompt plus generated tokens) and the raw prompt attention mass; mean over the last quarter of spans minus the first quarter.
\item \texttt{jmp}: on the 64-token block graph (each block's attention on earlier blocks, normalised per block), the mean attention-weighted backward distance in blocks.
\item $H_b$: entropy of $P(\cdot\mid p,t_{1:b},\sigma)$.
\item seg-25\% direction: $\langle \bar h^{(5)}_{[0.25L\pm128]}, d\rangle$, with $\bar h^{(5)}$ the mean 5th-layer state over the 256 tokens around position $0.25L$, and $d$ the normalised difference of class means on our labelled unfinished Olmo traces.
\item Min-K\% ($k=20$): mean of the lowest 20\% token log-probabilities of the problem text under the model.
\end{itemize}

\clearpage
\section{Builds scored on the leaderboard}
\label{app:ranking}
\begin{table}[H]
\centering\small
\caption{Leaderboard accuracy of our builds with their \textsc{Main} correct calls (each build's own recipe replayed on our \textsc{Main} traces). Per-model builds, the family \sone{} and \stwo{} belong to, score highest on both.}
\begin{tabular}{lcc}
\toprule
Build & Leaderboard & \textsc{Main} (43) \\
\midrule
v8-ent4096 (\sone{}) & 1.00 & 33 \\
v8-ent4096-olmoseg25 (\stwo{}) & 0.93 & 36 \\
v8 (base recipe, no single-feature readouts) & 0.93 & 32 \\
v7, v5 (per-model, router added in v7) & 0.93, 0.86 & 32, 34 \\
v6 (per-model, before the Olmo router) & 0.93 & 34 \\
max (per-model) & 0.86 & 33 \\
lookback (per-model) & 0.71 & 34 \\
cascade, v4, ensemble (per-model) & 0.64 & 30--34 \\
single-family builds (24) & 0.28--0.71 & 19--34 \\
\bottomrule
\end{tabular}
\end{table}

\section{Reproducibility}
\label{app:repro}
Code, recipes (\texttt{recipe.json}), shipped artifacts and the analysis scripts behind every number are at \todo{GitHub URL}: \texttt{submission\_core/} (inference), \texttt{analysis/build\_submission\_artifacts.py} (calibration), \texttt{analysis/report\_ablation.py} (Tables~\ref{tab:main}, \ref{tab:ablation} and \ref{tab:wb}) and \texttt{analysis/exp109\_mainv2.py} (Table~\ref{tab:position}).
Feature extraction and labelling ran on one NVIDIA RTX PRO 6000 and one A10G, using vLLM \citep{kwon2023vllm} for labelling samples and HF Transformers for features.
Model licences: Qwen3.5-4B, Skywork-OR1-Math-7B and Olmo-3-7B-Think are released under Apache-2.0, DeepSeek-R1-0528-Qwen3-8B under MIT \todo{verify on each model card}. The challenge data are used under the competition's terms and are not redistributed.

\clearpage
\section*{NeurIPS Paper Checklist}

\begin{enumerate}

\item {\bf Claims}
    \item[] Question: Do the main claims made in the abstract and introduction accurately reflect the paper's contributions and scope?
    \item[] Answer: \answerYes{} % Replace by \answerYes{}, \answerNo{}, or \answerNA{}.
    \item[] Justification: The abstract and Section~\ref{sec:intro} state results only on the official labels, with in-sample caveats (Section~\ref{sec:results} and the Limitations section).
    \item[] Guidelines:
    \begin{itemize}
        \item The answer \answerNA{} means that the abstract and introduction do not include the claims made in the paper.
        \item The abstract and/or introduction should clearly state the claims made, including the contributions made in the paper and important assumptions and limitations. A \answerNo{} or \answerNA{} answer to this question will not be perceived well by the reviewers. 
        \item The claims made should match theoretical and experimental results, and reflect how much the results can be expected to generalize to other settings. 
        \item It is fine to include aspirational goals as motivation as long as it is clear that these goals are not attained by the paper. 
    \end{itemize}

\item {\bf Limitations}
    \item[] Question: Does the paper discuss the limitations of the work performed by the authors?
    \item[] Answer: \answerYes{} % Replace by \answerYes{}, \answerNo{}, or \answerNA{}.
    \item[] Justification: See the Limitations section after the main text.
    \item[] Guidelines:
    \begin{itemize}
        \item The answer \answerNA{} means that the paper has no limitation while the answer \answerNo{} means that the paper has limitations, but those are not discussed in the paper. 
        \item The authors are encouraged to create a separate ``Limitations'' section in their paper.
        \item The paper should point out any strong assumptions and how robust the results are to violations of these assumptions (e.g., independence assumptions, noiseless settings, model well-specification, asymptotic approximations only holding locally). The authors should reflect on how these assumptions might be violated in practice and what the implications would be.
        \item The authors should reflect on the scope of the claims made, e.g., if the approach was only tested on a few datasets or with a few runs. In general, empirical results often depend on implicit assumptions, which should be articulated.
        \item The authors should reflect on the factors that influence the performance of the approach. For example, a facial recognition algorithm may perform poorly when image resolution is low or images are taken in low lighting. Or a speech-to-text system might not be used reliably to provide closed captions for online lectures because it fails to handle technical jargon.
        \item The authors should discuss the computational efficiency of the proposed algorithms and how they scale with dataset size.
        \item If applicable, the authors should discuss possible limitations of their approach to address problems of privacy and fairness.
        \item While the authors might fear that complete honesty about limitations might be used by reviewers as grounds for rejection, a worse outcome might be that reviewers discover limitations that aren't acknowledged in the paper. The authors should use their best judgment and recognize that individual actions in favor of transparency play an important role in developing norms that preserve the integrity of the community. Reviewers will be specifically instructed to not penalize honesty concerning limitations.
    \end{itemize}

\item {\bf Theory assumptions and proofs}
    \item[] Question: For each theoretical result, does the paper provide the full set of assumptions and a complete (and correct) proof?
    \item[] Answer: \answerNA{} % Replace by \answerYes{}, \answerNo{}, or \answerNA{}.
    \item[] Justification: The paper has no theoretical results.
    \item[] Guidelines:
    \begin{itemize}
        \item The answer \answerNA{} means that the paper does not include theoretical results. 
        \item All the theorems, formulas, and proofs in the paper should be numbered and cross-referenced.
        \item All assumptions should be clearly stated or referenced in the statement of any theorems.
        \item The proofs can either appear in the main paper or the supplemental material, but if they appear in the supplemental material, the authors are encouraged to provide a short proof sketch to provide intuition. 
        \item Inversely, any informal proof provided in the core of the paper should be complemented by formal proofs provided in appendix or supplemental material.
        \item Theorems and Lemmas that the proof relies upon should be properly referenced. 
    \end{itemize}

    \item {\bf Experimental result reproducibility}
    \item[] Question: Does the paper fully disclose all the information needed to reproduce the main experimental results of the paper to the extent that it affects the main claims and/or conclusions of the paper (regardless of whether the code and data are provided or not)?
    \item[] Answer: \answerYes{} % Replace by \answerYes{}, \answerNo{}, or \answerNA{}.
    \item[] Justification: Section~\ref{sec:method} and Appendix~\ref{app:features} define every feature, threshold and decision rule; the shipped \texttt{recipe.json} files hold all values.
    \item[] Guidelines:
    \begin{itemize}
        \item The answer \answerNA{} means that the paper does not include experiments.
        \item If the paper includes experiments, a \answerNo{} answer to this question will not be perceived well by the reviewers: Making the paper reproducible is important, regardless of whether the code and data are provided or not.
        \item If the contribution is a dataset and\slash or model, the authors should describe the steps taken to make their results reproducible or verifiable. 
        \item Depending on the contribution, reproducibility can be accomplished in various ways. For example, if the contribution is a novel architecture, describing the architecture fully might suffice, or if the contribution is a specific model and empirical evaluation, it may be necessary to either make it possible for others to replicate the model with the same dataset, or provide access to the model. In general. releasing code and data is often one good way to accomplish this, but reproducibility can also be provided via detailed instructions for how to replicate the results, access to a hosted model (e.g., in the case of a large language model), releasing of a model checkpoint, or other means that are appropriate to the research performed.
        \item While NeurIPS does not require releasing code, the conference does require all submissions to provide some reasonable avenue for reproducibility, which may depend on the nature of the contribution. For example
        \begin{enumerate}
            \item If the contribution is primarily a new algorithm, the paper should make it clear how to reproduce that algorithm.
            \item If the contribution is primarily a new model architecture, the paper should describe the architecture clearly and fully.
            \item If the contribution is a new model (e.g., a large language model), then there should either be a way to access this model for reproducing the results or a way to reproduce the model (e.g., with an open-source dataset or instructions for how to construct the dataset).
            \item We recognize that reproducibility may be tricky in some cases, in which case authors are welcome to describe the particular way they provide for reproducibility. In the case of closed-source models, it may be that access to the model is limited in some way (e.g., to registered users), but it should be possible for other researchers to have some path to reproducing or verifying the results.
        \end{enumerate}
    \end{itemize}


\item {\bf Open access to data and code}
    \item[] Question: Does the paper provide open access to the data and code, with sufficient instructions to faithfully reproduce the main experimental results, as described in supplemental material?
    \item[] Answer: \answerYes{} % Replace by \answerYes{}, \answerNo{}, or \answerNA{}.
    \item[] Justification: Code, recipes and analysis scripts are released (Appendix~\ref{app:repro}); the official challenge data are not ours to redistribute.
    \item[] Guidelines:
    \begin{itemize}
        \item The answer \answerNA{} means that paper does not include experiments requiring code.
        \item Please see the NeurIPS code and data submission guidelines (\url{https://neurips.cc/public/guides/CodeSubmissionPolicy}) for more details.
        \item While we encourage the release of code and data, we understand that this might not be possible, so \answerNo{} is an acceptable answer. Papers cannot be rejected simply for not including code, unless this is central to the contribution (e.g., for a new open-source benchmark).
        \item The instructions should contain the exact command and environment needed to run to reproduce the results. See the NeurIPS code and data submission guidelines (\url{https://neurips.cc/public/guides/CodeSubmissionPolicy}) for more details.
        \item The authors should provide instructions on data access and preparation, including how to access the raw data, preprocessed data, intermediate data, and generated data, etc.
        \item The authors should provide scripts to reproduce all experimental results for the new proposed method and baselines. If only a subset of experiments are reproducible, they should state which ones are omitted from the script and why.
        \item At submission time, to preserve anonymity, the authors should release anonymized versions (if applicable).
        \item Providing as much information as possible in supplemental material (appended to the paper) is recommended, but including URLs to data and code is permitted.
    \end{itemize}


\item {\bf Experimental setting/details}
    \item[] Question: Does the paper specify all the training and test details (e.g., data splits, hyperparameters, how they were chosen, type of optimizer) necessary to understand the results?
    \item[] Answer: \answerYes{} % Replace by \answerYes{}, \answerNo{}, or \answerNA{}.
    \item[] Justification: Sections~\ref{sec:data}--\ref{sec:method}: data splits, how signs, thresholds and recipes were chosen; nothing is trained by gradient descent.
    \item[] Guidelines:
    \begin{itemize}
        \item The answer \answerNA{} means that the paper does not include experiments.
        \item The experimental setting should be presented in the core of the paper to a level of detail that is necessary to appreciate the results and make sense of them.
        \item The full details can be provided either with the code, in appendix, or as supplemental material.
    \end{itemize}

\item {\bf Experiment statistical significance}
    \item[] Question: Does the paper report error bars suitably and correctly defined or other appropriate information about the statistical significance of the experiments?
    \item[] Answer: \answerYes{} % Replace by \answerYes{}, \answerNo{}, or \answerNA{}.
    \item[] Justification: Table~\ref{tab:main} gives Wilson 95\% intervals; Section~\ref{sec:results} and the Limitations note that the two submissions do not differ significantly.
    \item[] Guidelines:
    \begin{itemize}
        \item The answer \answerNA{} means that the paper does not include experiments.
        \item The authors should answer \answerYes{} if the results are accompanied by error bars, confidence intervals, or statistical significance tests, at least for the experiments that support the main claims of the paper.
        \item The factors of variability that the error bars are capturing should be clearly stated (for example, train/test split, initialization, random drawing of some parameter, or overall run with given experimental conditions).
        \item The method for calculating the error bars should be explained (closed form formula, call to a library function, bootstrap, etc.)
        \item The assumptions made should be given (e.g., Normally distributed errors).
        \item It should be clear whether the error bar is the standard deviation or the standard error of the mean.
        \item It is OK to report 1-sigma error bars, but one should state it. The authors should preferably report a 2-sigma error bar than state that they have a 96\% CI, if the hypothesis of Normality of errors is not verified.
        \item For asymmetric distributions, the authors should be careful not to show in tables or figures symmetric error bars that would yield results that are out of range (e.g., negative error rates).
        \item If error bars are reported in tables or plots, the authors should explain in the text how they were calculated and reference the corresponding figures or tables in the text.
    \end{itemize}

\item {\bf Experiments compute resources}
    \item[] Question: For each experiment, does the paper provide sufficient information on the computer resources (type of compute workers, memory, time of execution) needed to reproduce the experiments?
    \item[] Answer: \answerYes{} % Replace by \answerYes{}, \answerNo{}, or \answerNA{}.
    \item[] Justification: Per-problem cost in Section~\ref{sec:method}; hardware in Appendix~\ref{app:repro}.
    \item[] Guidelines:
    \begin{itemize}
        \item The answer \answerNA{} means that the paper does not include experiments.
        \item The paper should indicate the type of compute workers CPU or GPU, internal cluster, or cloud provider, including relevant memory and storage.
        \item The paper should provide the amount of compute required for each of the individual experimental runs as well as estimate the total compute. 
        \item The paper should disclose whether the full research project required more compute than the experiments reported in the paper (e.g., preliminary or failed experiments that didn't make it into the paper). 
    \end{itemize}
    
\item {\bf Code of ethics}
    \item[] Question: Does the research conducted in the paper conform, in every respect, with the NeurIPS Code of Ethics \url{https://neurips.cc/public/EthicsGuidelines}?
    \item[] Answer: \answerYes{} % Replace by \answerYes{}, \answerNo{}, or \answerNA{}.
    \item[] Justification: The work uses public models and challenge data and involves no human subjects.
    \item[] Guidelines:
    \begin{itemize}
        \item The answer \answerNA{} means that the authors have not reviewed the NeurIPS Code of Ethics.
        \item If the authors answer \answerNo, they should explain the special circumstances that require a deviation from the Code of Ethics.
        \item The authors should make sure to preserve anonymity (e.g., if there is a special consideration due to laws or regulations in their jurisdiction).
    \end{itemize}


\item {\bf Broader impacts}
    \item[] Question: Does the paper discuss both potential positive societal impacts and negative societal impacts of the work performed?
    \item[] Answer: \answerNA{} % Replace by \answerYes{}, \answerNo{}, or \answerNA{}.
    \item[] Justification: The work is a diagnostic study of public models' reasoning robustness with no direct deployment path or foreseeable negative societal impact.
    \item[] Guidelines:
    \begin{itemize}
        \item The answer \answerNA{} means that there is no societal impact of the work performed.
        \item If the authors answer \answerNA{} or \answerNo, they should explain why their work has no societal impact or why the paper does not address societal impact.
        \item Examples of negative societal impacts include potential malicious or unintended uses (e.g., disinformation, generating fake profiles, surveillance), fairness considerations (e.g., deployment of technologies that could make decisions that unfairly impact specific groups), privacy considerations, and security considerations.
        \item The conference expects that many papers will be foundational research and not tied to particular applications, let alone deployments. However, if there is a direct path to any negative applications, the authors should point it out. For example, it is legitimate to point out that an improvement in the quality of generative models could be used to generate Deepfakes for disinformation. On the other hand, it is not needed to point out that a generic algorithm for optimizing neural networks could enable people to train models that generate Deepfakes faster.
        \item The authors should consider possible harms that could arise when the technology is being used as intended and functioning correctly, harms that could arise when the technology is being used as intended but gives incorrect results, and harms following from (intentional or unintentional) misuse of the technology.
        \item If there are negative societal impacts, the authors could also discuss possible mitigation strategies (e.g., gated release of models, providing defenses in addition to attacks, mechanisms for monitoring misuse, mechanisms to monitor how a system learns from feedback over time, improving the efficiency and accessibility of ML).
    \end{itemize}
    
\item {\bf Safeguards}
    \item[] Question: Does the paper describe safeguards that have been put in place for responsible release of data or models that have a high risk for misuse (e.g., pre-trained language models, image generators, or scraped datasets)?
    \item[] Answer: \answerNA{} % Replace by \answerYes{}, \answerNo{}, or \answerNA{}.
    \item[] Justification: We release no data or models with a high risk of misuse.
    \item[] Guidelines:
    \begin{itemize}
        \item The answer \answerNA{} means that the paper poses no such risks.
        \item Released models that have a high risk for misuse or dual-use should be released with necessary safeguards to allow for controlled use of the model, for example by requiring that users adhere to usage guidelines or restrictions to access the model or implementing safety filters. 
        \item Datasets that have been scraped from the Internet could pose safety risks. The authors should describe how they avoided releasing unsafe images.
        \item We recognize that providing effective safeguards is challenging, and many papers do not require this, but we encourage authors to take this into account and make a best faith effort.
    \end{itemize}

\item {\bf Licenses for existing assets}
    \item[] Question: Are the creators or original owners of assets (e.g., code, data, models), used in the paper, properly credited and are the license and terms of use explicitly mentioned and properly respected?
    \item[] Answer: \answerYes{} % Replace by \answerYes{}, \answerNo{}, or \answerNA{}.
    \item[] Justification: Models, libraries and methods are cited; licences are listed in Appendix~\ref{app:repro}.
    \item[] Guidelines:
    \begin{itemize}
        \item The answer \answerNA{} means that the paper does not use existing assets.
        \item The authors should cite the original paper that produced the code package or dataset.
        \item The authors should state which version of the asset is used and, if possible, include a URL.
        \item The name of the license (e.g., CC-BY 4.0) should be included for each asset.
        \item For scraped data from a particular source (e.g., website), the copyright and terms of service of that source should be provided.
        \item If assets are released, the license, copyright information, and terms of use in the package should be provided. For popular datasets, \url{paperswithcode.com/datasets} has curated licenses for some datasets. Their licensing guide can help determine the license of a dataset.
        \item For existing datasets that are re-packaged, both the original license and the license of the derived asset (if it has changed) should be provided.
        \item If this information is not available online, the authors are encouraged to reach out to the asset's creators.
    \end{itemize}

\item {\bf New assets}
    \item[] Question: Are new assets introduced in the paper well documented and is the documentation provided alongside the assets?
    \item[] Answer: \answerYes{} % Replace by \answerYes{}, \answerNo{}, or \answerNA{}.
    \item[] Justification: The released code and recipe artifacts are documented in the repository (Appendix~\ref{app:repro}).
    \item[] Guidelines:
    \begin{itemize}
        \item The answer \answerNA{} means that the paper does not release new assets.
        \item Researchers should communicate the details of the dataset\slash code\slash model as part of their submissions via structured templates. This includes details about training, license, limitations, etc. 
        \item The paper should discuss whether and how consent was obtained from people whose asset is used.
        \item At submission time, remember to anonymize your assets (if applicable). You can either create an anonymized URL or include an anonymized zip file.
    \end{itemize}

\item {\bf Crowdsourcing and research with human subjects}
    \item[] Question: For crowdsourcing experiments and research with human subjects, does the paper include the full text of instructions given to participants and screenshots, if applicable, as well as details about compensation (if any)? 
    \item[] Answer: \answerNA{} % Replace by \answerYes{}, \answerNo{}, or \answerNA{}.
    \item[] Justification: No crowdsourcing or human-subject research.
    \item[] Guidelines:
    \begin{itemize}
        \item The answer \answerNA{} means that the paper does not involve crowdsourcing nor research with human subjects.
        \item Including this information in the supplemental material is fine, but if the main contribution of the paper involves human subjects, then as much detail as possible should be included in the main paper. 
        \item According to the NeurIPS Code of Ethics, workers involved in data collection, curation, or other labor should be paid at least the minimum wage in the country of the data collector. 
    \end{itemize}

\item {\bf Institutional review board (IRB) approvals or equivalent for research with human subjects}
    \item[] Question: Does the paper describe potential risks incurred by study participants, whether such risks were disclosed to the subjects, and whether Institutional Review Board (IRB) approvals (or an equivalent approval/review based on the requirements of your country or institution) were obtained?
    \item[] Answer: \answerNA{} % Replace by \answerYes{}, \answerNo{}, or \answerNA{}.
    \item[] Justification: No human subjects.
    \item[] Guidelines:
    \begin{itemize}
        \item The answer \answerNA{} means that the paper does not involve crowdsourcing nor research with human subjects.
        \item Depending on the country in which research is conducted, IRB approval (or equivalent) may be required for any human subjects research. If you obtained IRB approval, you should clearly state this in the paper. 
        \item We recognize that the procedures for this may vary significantly between institutions and locations, and we expect authors to adhere to the NeurIPS Code of Ethics and the guidelines for their institution. 
        \item For initial submissions, do not include any information that would break anonymity (if applicable), such as the institution conducting the review.
    \end{itemize}

\item {\bf Declaration of LLM usage}
    \item[] Question: Does the paper describe the usage of LLMs if it is an important, original, or non-standard component of the core methods in this research? Note that if the LLM is used only for writing, editing, or formatting purposes and does \emph{not} impact the core methodology, scientific rigor, or originality of the research, declaration is not required.
    %this research? 
    \item[] Answer: \answerYes{} % Replace by \answerYes{}, \answerNo{}, or \answerNA{}.
    \item[] Justification: An LLM (Claude Sonnet) wrote the perturbation texts for our own calibration labels, and a second open model (Qwen) acts as a judge for Olmo (Section~\ref{sec:method}); all LLM use is described in the AI usage statement.
    \item[] Guidelines:
    \begin{itemize}
        \item The answer \answerNA{} means that the core method development in this research does not involve LLMs as any important, original, or non-standard components.
        \item Please refer to our LLM policy in the NeurIPS handbook for what should or should not be described.
    \end{itemize}

\end{enumerate}

\end{document}
